\section{Introduction}
\label{sec:intro}

%% Leave page number of the first page empty
%% 
\thispagestyle{empty}

This is the template file for writing your bachelor’s, master’s or licentiate 
thesis. The template contains parts and text to accommodate the different degree 
levels, bachelor’s and master’s, and so it may contain parts that may not be 
relevant to your thesis. If so, simply delete those parts in the thesis that you
do not need. This applies particularly to the abstract pages, the list of
symbols, the list of abbreviations, and appendices. You may also want to adjust
the template section titles to better suit your work.

The text and comments in this template advise you how to use the template, give 
general pointers on how to write your thesis (you must still refer to writing
guides and converse with your advisor and/or supervisor in this regard), and
provide the technical specifications for the layout and style of the template.

\subsection{Typical content in the introduction}

In principle, the introduction is like the abstract, only broader in scope, and
more detailed. The introduction generally describes the following:

%% An example of an unordered (itemised) list. Instead on the default bullet
%% point, an endash (-- in LaTeX code) is used as the item label.
\begin{itemize}
	\item[--] a description of the background of the field of study, what
	similar work others have already done, as well as an overview of the study,
	\item[--] the goals of the study,
	\item[--] the primary research question and the sub-problems in the line
	of inquiry, and
	\item[--] the scope and constraints of the study along with the main
	concepts involved.
\end{itemize}

Although the introduction is a general description of the study, be concise and 
avoid writing a lengthy introduction. A concise introduction need not have any 
subsections.

\clearpage
